% Generated draft: requires mathematical and source review.
\documentclass[11pt]{article}
\usepackage[margin=25mm]{geometry}
\usepackage{fontspec}
\setmainfont{Latin Modern Roman}
\usepackage{amsmath,amssymb,mathtools,unicode-math}
\setmathfont{Latin Modern Math}
\usepackage{enumitem,needspace,xurl,hyperref}
\setlength{\parindent}{0pt}
\setlength{\parskip}{5pt}
\setlength{\emergencystretch}{3em}
\newcommand{\N}{\mathbb{N}}
\newcommand{\Z}{\mathbb{Z}}
\newcommand{\Q}{\mathbb{Q}}
\newcommand{\R}{\mathbb{R}}
\newcommand{\C}{\mathbb{C}}
\newcommand{\F}{\mathbb{F}}
\newcommand{\A}{\mathbb{A}}
\newcommand{\PP}{\mathbb{P}}
\newcommand{\E}{\mathbb{E}}
\newcommand{\eps}{\varepsilon}
\newcommand{\dd}{\,\mathrm{d}}
\begin{document}
\section*{UnsolvedMath ID: 30005436}
\textbf{Source problem code:} catalogue entry
\subsection{Definitions and mathematical statement}
Let \(K\) be a field and let \(E=K^n\). Let \(A\in M_n(K)\) define the \(K\)-bilinear form
\[
 b(v,w)=v^{\mathsf T}Aw.
\]
Assume that \(b\) is anisotropic, meaning that
\[
 b(v,v)\ne 0\qquad\text{for every }v\in E\setminus\{0\}.
\]
Symmetry of \(A\) is not assumed. Let \(L(K,n)\) denote the lattice of linear subspaces of \(E\). For a subspace \(U\), use the right orthogonal complement
\[
 U^{\perp_b}=\{z\in E: b(u,z)=0\text{ for every }u\in U\}.
\]
The source associates to \(b\) an \(L\)-algebra \(X(b)\) on this lattice by the Sasaki operation
\[
 U\mathbin{\cdot}W=U^{\perp_b}+(U\cap W),
 \qquad U,W\in L(K,n),
\]
where \(+\) is the sum of subspaces. Here the orthogonal complement is interpreted in the ambient space \(E\): the extracted source instead contains an unclear membership condition involving \(V\). This interpretation requires checking against the original notation.

Write \(G(X(b))\) for the structure group attached to this \(L\)-algebra in the source's convention. The supplied excerpt does not give its defining presentation or the general definition of an \(L\)-algebra, so these conventions remain to be supplied for a fully self-contained formulation.

For comparison, when \(A=A^{\mathsf T}\), the supplied excerpt states, as Theorem 7 with reference label [1], that
\[
 G(X(b))\cong\operatorname{PPU}(b),
\]
where
\[
 \operatorname{PPU}(b)=
 \left\{\varphi(t)\in\operatorname{GL}_n\bigl(K[t,t^{-1}]\bigr):
 \begin{aligned}
 \varphi(t^{-1})^{\mathsf T}A\varphi(t)&=A,\\
 \varphi(1)&=I_n
 \end{aligned}
 \right\}.
\]
Here \(t\) is a formal indeterminate, \(K[t,t^{-1}]\) is the Laurent polynomial ring, and the displayed transpose and substitution do not introduce an involution on \(K\).

\textbf{Problem 8.} For an arbitrary anisotropic \(K\)-bilinear form \(b\) as above, with \(A\) not necessarily symmetric, find a representation or concrete description of \(G(X(b))\) analogous to the displayed pure paraunitary description. The source does not specify the class of admissible representations or a precise criterion for being analogous; in particular, it does not assert that the same displayed matrix group describes the nonsymmetric case.

\textbf{Unresolved contextual qualification.} Immediately following Problem 8, the extraction lists: (1) \(n\in\{1,2\}\), marked ``trivial''; (2) existence of an anisotropic hermitian bilinear form on \(K^n\), with reference label [3]; and (3) \(K\) a finite field, with reference label [2]. It then states that a Desarguesian lattice without a brick structure must be infinite. The supplied text does not identify the assertion governed by this list, define ``brick structure'', or specify the involution intended in the hermitian condition. These fragments are retained as unresolved context and are not treated as additional hypotheses or established special cases of Problem 8.
\subsection{Short English statement}
Given an anisotropic bilinear form over a field, describe the structure group of its associated subspace \(L\)-algebra when the form may be nonsymmetric. The requested description should be analogous to the pure paraunitary model supplied for symmetric forms. The structure-group convention, the precise meaning of an analogous representation, and several extracted contextual qualifications require human review.
\subsection{Sources}
\begin{itemize}
\item Local UnsolvedMath registry: \texttt{data/unsolvedmath/problems.json}, permanent ID \(30005436\), title ``Structure Groups of Anisotropic Bilinear Forms'', source code \texttt{catalogue entry}. The supplied record identifies the target as Problem 8 in the source excerpt and gives proposed year 2023; no individual proposer is supplied.
\item Supplied source citation: \emph{Mini-Workshop: Skew Braces and the Yang--Baxter Equation} (2023), \emph{Oberwolfach Reports}, pp.\ 537--563, DOI \texttt{10.4171/owr/2023/9}. Source URL: \url{https://doi.org/10.4171/owr/2023/9}.
\item Supplied literature source, labelled ``original OWR report'': \url{https://doi.org/10.4171/owr/2023/9}.
\item Literature assessment dated \texttt{2026-08-22}: the supplied curation labels the problem open and reports that no authoritative later solution was located. This is supplied dated curation, not an independently verified literature assessment.
\item The registry labels the statement \texttt{corrected\_verified} and reports a direct check of the official report. That verification claim is part of the supplied record and has not been independently checked for this draft. The excerpt's reference labels [1], [2], and [3] have no accompanying bibliographic entries in the supplied material.
\end{itemize}
\end{document}
